% Referencia estética elegida por el usuario: TesisPrueba.pdf (2026-10-29).
% Medidas y adaptaciones documentadas en estilo.md; contenido según tipo 1.
\usepackage{fontspec}
\setmainfont{Latin Modern Roman}
\setsansfont{Latin Modern Sans}
% Unicode directo; evitar la incompatibilidad de shorthands con los hooks APA.
\usepackage[spanish,es-tabla,es-noshorthands,es-lcroman]{babel}
\usepackage[margin=1in,headheight=15pt,headsep=18pt,footskip=30pt]{geometry}
\usepackage{setspace}
\onehalfspacing
\setlength{\parindent}{1.5em}
\setlength{\parskip}{0pt}
\usepackage{titlesec}
\titleformat{\chapter}[display]
  {\normalfont\Huge\bfseries}{\chaptertitlename\ \thechapter}{20pt}{}
\titlespacing*{\chapter}{0pt}{50pt}{40pt}
\titleformat{\section}{\normalfont\Large\bfseries}{\thesection.}{1em}{}
\titleformat{\subsection}{\normalfont\large\bfseries}{\thesubsection.}{1em}{}
\usepackage{titletoc}
% Reservar espacio también para capítulos de dos dígitos en el índice.
\titlecontents{chapter}[2.5em]{\addvspace{1em}\bfseries}
  {\contentslabel{2.5em}}{\hspace*{-2.5em}}{\hfill\contentspage}
\usepackage{fancyhdr}
\pagestyle{fancy}
\fancyhf{}
\fancyhead[LE,RO]{\small\bfseries\thepage}
\fancyhead[LO]{\small\sffamily\nouppercase{\rightmark}}
\fancyhead[RE]{\small\sffamily\bfseries\nouppercase{\leftmark}}
\renewcommand{\headrulewidth}{0.4pt}
\renewcommand{\footrulewidth}{0pt}
\renewcommand{\chaptermark}[1]{\markboth{#1}{}}
\renewcommand{\sectionmark}[1]{\markright{\thesection\ #1}}
\fancypagestyle{plain}{%
  \fancyhf{}%
  \fancyfoot[C]{\thepage}%
  \renewcommand{\headrulewidth}{0pt}%
  \renewcommand{\footrulewidth}{0pt}%
}
\usepackage{graphicx}
\usepackage{booktabs}
\usepackage{longtable}
\usepackage{csquotes}
\usepackage[backend=biber,style=apa]{biblatex}
\addbibresource{bibliografia/referencias.bib}
\usepackage[hidelinks]{hyperref}
\hypersetup{pdftitle={Trabajo Final -- borrador},pdfsubject={Tipo 1: Proyecto y/o Diseño de Ingeniería}}
\setcounter{tocdepth}{2}
\setcounter{secnumdepth}{2}
% Los pendientes son visibles también en el PDF para evitar falsas completitudes.
\newcommand{\pendiente}[1]{\par\noindent\textbf{[COMPLETAR]} #1\par}
